\section{A tétel lényege és vizsgaszintű vázlata}

\begin{summarybox}
A rendezés olyan alapvető algoritmikai feladat, amelyben egy sorozat elemeit egy megadott rendezési reláció szerint kell sorrendbe állítani. A tétel fő gondolata az, hogy a rendezési algoritmusok között nagy különbség van aszerint, hogy összehasonlításokon alapulnak-e, helyben dolgoznak-e, stabilak-e, mennyi segédmemóriát igényelnek, és milyen időigényük van a legjobb, átlagos és legrosszabb esetben.
\end{summarybox}

Vizsgán a tétel jól felépíthető a következő sorrendben:

\begin{enumerate}
  \item rendezés fogalma, rendezési reláció, kulcs és rekord;
  \item rendezési szempontok: stabilitás, helyben rendezés, adaptivitás, belső és külső rendezés;
  \item egyszerű, többnyire négyzetes rendezések: beszúró, buborék, minimum kiválasztásos;
  \item oszd meg és uralkodj elv: összefésülő rendezés és gyorsrendezés;
  \item Shell rendezés és négyzetes rendezés mint köztes ötletek;
  \item lineáris idejű rendezések: leszámláló és számjegyes rendezés;
  \item összehasonlító rendezések alsó korlátja: bármely összehasonlító rendezés legrosszabb esetben legalább $\Omega(n\log n)$ összehasonlítást igényel.
\end{enumerate}

A legfontosabb vizsgamondat: az összehasonlító rendezések általános alsó korlátja $\Omega(n\log n)$, ezért lineáris idejű rendezést csak akkor kaphatunk, ha a kulcsokra külön feltételeket teszünk, és nem pusztán páronkénti összehasonlításokból dolgozunk.

\begin{center}
\begingroup
\setlength{\tabcolsep}{4pt}
\begin{longtable}{L{0.20\textwidth}L{0.34\textwidth}L{0.23\textwidth}L{0.13\textwidth}}
\toprule
Rendezés & Alapelv & Tipikus tulajdonság & Időigény röviden \\
\midrule
\endfirsthead
\toprule
Rendezés & Alapelv & Tipikus tulajdonság & Időigény röviden \\
\midrule
\endhead
Beszúró rendezés & A következő elemet beszúrja az előtte lévő rendezett részbe. & Stabil, helyben rendez, majdnem rendezett bemeneten jó. & legjobb $\Theta(n)$, rossz $\Theta(n^2)$ \\
Összefésülő rendezés & Két félre bont, rekurzívan rendez, majd összefésül. & Stabil lehet, de extra tárat igényel. & $\Theta(n\log n)$ \\
Gyorsrendezés & Pivot szerint particionál, majd rekurzívan rendezi a részeket. & Gyakorlatban gyors, helyben rendezhető, nem stabil. & átlag $\Theta(n\log n)$, rossz $\Theta(n^2)$ \\
Buborék rendezés & Szomszédos elemeket cserél, amíg van rossz sorrend. & Egyszerű, sok cserét végez. & legjobb $\Theta(n)$, rossz $\Theta(n^2)$ \\
Shell rendezés & Fogyó növekményekkel távoli elemeket rendez beszúró jelleggel. & Időigénye a növekményektől függ. & gyakran kb. $\Theta(n^{1.5})$ \\
Minimum kiválasztásos rendezés & Mindig kiválasztja a maradék minimumát és előre cseréli. & Kevés csere, de sok összehasonlítás. & $\Theta(n^2)$ \\
Négyzetes rendezés & Csoportokra bontva gyorsítja a minimumkiválasztás gondolatát. & Köztes módszer, csoportos minimumkezeléssel. & kb. $\Theta(n^{1.5})$ \\
Leszámláló rendezés & Megszámolja a kulcsok előfordulását és prefixösszegekkel pozíciót számít. & Nem összehasonlító, egész kulcstartomány kell. & $\Theta(n+k)$ \\
Számjegyes rendezés & Több stabil részrendezést végez számjegyenként vagy pozíciónként. & Nem összehasonlító, stabil segédrendezés kell. & $\Theta(d(n+k))$ \\
\bottomrule
\end{longtable}
\endgroup

\end{center}

\section{A rendezés alapfogalmai}

\subsection{Rendezési reláció és rendezendő sorozat}

Rendezésről akkor beszélhetünk, ha az adatok között értelmezve van valamilyen rendezési reláció. Tipikus esetben a reláció a $\leq$ reláció, de a rendezés alapja lehet például név szerinti betűrend, dátum, ár, pontszám vagy többmezős kulcs is.

Egy $A[1..n]$ sorozat növekvő rendezése azt jelenti, hogy a kimeneti sorozat ugyanazokat az elemeket tartalmazza, mint a bemenet, de teljesül rá:

\[
A[1] \leq A[2] \leq \dots \leq A[n].
\]

Ha rekordokat rendezünk, akkor általában nem az egész rekord határozza meg a sorrendet, hanem annak egy vagy több \textbf{kulcsa}. Például hallgatói rekordoknál lehet kulcs a Neptun-kód, a név vagy az átlag.

\subsection{Rendezési algoritmusok fontos tulajdonságai}

Egy rendezési algoritmust nem csak az időigény alapján érdemes értékelni. A gyakorlatban a következő tulajdonságok is fontosak:

\begin{itemize}
  \item \textbf{Helyben rendezés:} az algoritmus legfeljebb konstans vagy nagyon kicsi extra tárat használ, az eredmény az eredeti tömbben áll elő.
  \item \textbf{Stabilitás:} az azonos kulcsú elemek relatív sorrendje megmarad.
  \item \textbf{Adaptivitás:} az algoritmus gyorsabb lehet, ha a bemenet már majdnem rendezett.
  \item \textbf{Belső rendezés:} az adatok elférnek a memóriában.
  \item \textbf{Külső rendezés:} az adatok háttértáron vannak, ezért a fájlműveletek és a blokkos hozzáférés dominálnak.
  \item \textbf{Összehasonlító rendezés:} a sorrend kizárólag elempárok összehasonlításából derül ki.
  \item \textbf{Nem összehasonlító rendezés:} a kulcs szerkezetét vagy tartományát is kihasználja, például egész kulcsok kis intervallumát.
\end{itemize}

\begin{center}
\begingroup
\setlength{\tabcolsep}{4pt}
\begin{tabularx}{\textwidth}{L{0.24\textwidth}X L{0.26\textwidth}}
\toprule
Tulajdonság & Jelentés & Tipikus példa \\
\midrule
Helyben rendezés & Az eredmény az eredeti tömbben áll elő, és az algoritmus csak konstans vagy kicsi extra tárat használ. & beszúró, gyorsrendezés, minimum kiválasztásos \\
Stabilitás & Az azonos kulcsú elemek megtartják egymáshoz képesti eredeti sorrendjüket. & beszúró, összefésülő, leszámláló \\
Adaptivitás & Az algoritmus kihasználja, ha a bemenet már részben rendezett. & beszúró, optimalizált buborék \\
Összehasonlító rendezés & A sorrend elempárok összehasonlításából derül ki. & beszúró, merge sort, quicksort \\
Nem összehasonlító rendezés & A kulcsok tartományát vagy szerkezetét is kihasználja. & leszámláló, számjegyes \\
Belső rendezés & Az adatok a főmemóriában vannak. & a legtöbb tömbös rendezés \\
Külső rendezés & Az adatok háttértáron vannak, ezért az I/O-költség is fontos. & külső összefésülés \\
\bottomrule
\end{tabularx}
\endgroup

\end{center}

\begin{center}
\begin{tikzpicture}[node distance=7mm, scale=0.82, transform shape]
  \node[flow, minimum width=34mm] (root) {Rendezési\\algoritmusok};
  \node[flow, below left=12mm and 8mm of root, minimum width=34mm] (comp) {Összehasonlító\\rendezések};
  \node[flow, below right=12mm and 8mm of root, minimum width=34mm] (noncomp) {Nem összehasonlító\\rendezések};
  \node[flow, below left=10mm and -3mm of comp, minimum width=29mm] (simple) {Egyszerű\\$\Theta(n^2)$\\módszerek};
  \node[flow, below right=10mm and -3mm of comp, minimum width=29mm] (divide) {Oszd meg\\és uralkodj};
  \node[flow, below=10mm of noncomp, minimum width=34mm] (linear) {Kulcstartományt\\kihasználó\\módszerek};
  \node[flow, below=8mm of simple, minimum width=32mm] (s1) {beszúró\\buborék\\minimum kiválasztásos};
  \node[flow, below=8mm of divide, minimum width=31mm] (s2) {összefésülő\\gyorsrendezés};
  \node[flow, below=8mm of linear, minimum width=31mm] (s3) {leszámláló\\számjegyes};
  \draw[arrow] (root) -- (comp);
  \draw[arrow] (root) -- (noncomp);
  \draw[arrow] (comp) -- (simple);
  \draw[arrow] (comp) -- (divide);
  \draw[arrow] (noncomp) -- (linear);
  \draw[arrow] (simple) -- (s1);
  \draw[arrow] (divide) -- (s2);
  \draw[arrow] (linear) -- (s3);
\end{tikzpicture}
\end{center}


\subsection{Időelemzés szempontjai}

A rendezések időigényét gyakran három esetben vizsgáljuk:

\begin{itemize}
  \item \textbf{legjobb eset:} a bemenet különösen kedvező, például már rendezett;
  \item \textbf{átlagos eset:} a bemenet tipikus vagy véletlenszerű;
  \item \textbf{legrosszabb eset:} a bemenet az algoritmus számára a lehető legkedvezőtlenebb.
\end{itemize}

A rendezések összehasonlításánál gyakori jelölések:

\[
\Theta(n^2),\qquad \Theta(n\log n),\qquad \Theta(n+k),\qquad \Theta(d(n+k)).
\]

Itt $n$ az elemek száma, $k$ például a kulcsok lehetséges értékeinek száma, $d$ pedig számjegyes rendezésnél a feldolgozott pozíciók száma.

\section{Beszúró rendezés}

\subsection{Alapelv}

A \textbf{beszúró rendezés} úgy működik, ahogyan kézzel is rendezünk például kártyákat. A sorozat elején mindig fenntart egy rendezett részt, majd a következő elemet beszúrja ebbe a rendezett részbe a megfelelő helyre.

A legfontosabb invariáns:

\begin{summarybox}
A $j$-edik külső cikluslépés előtt az $A[1..j-1]$ rész már rendezett, és pontosan az eredeti első $j-1$ elemet tartalmazza rendezett sorrendben.
\end{summarybox}

\begin{center}
\begin{tikzpicture}[x=0.78cm,y=0.78cm]
  \foreach \i/\val in {0/1,1/3,2/5,3/8,4/4,5/7,6/2,7/9} {
    \draw[fill=lightaccent] (\i,0) rectangle +(1,0.8);
    \node at (\i+0.5,0.4) {$\val$};
  }
  \draw[very thick, accent] (0,-0.18) -- (4,-0.18);
  \node[align=center, font=\small] at (2,-0.78) {rendezett rész:\\$A[1..j-1]$};
  \draw[very thick, orange!80!black] (4,-0.18) -- (5,-0.18);
  \node[align=center, font=\small] at (4.5,-1.10) {beszúrandó\\kulcs};
  \draw[->, thick, orange!80!black] (4.5,1.50) .. controls (4.0,1.05) and (3.4,0.95) .. (2.7,0.86);
  \node[align=center, font=\small] at (4.7,2.05) {a kulcs balra tolódik,\\amíg megvan a helye};
\end{tikzpicture}
\end{center}


\subsection{Pszeudókód}

\zvpseudocode[caption={Beszúró rendezés},label={lst:zv05_beszuro_rendezes}]{src/code/zv05_beszuro_rendezes.pseudo}

\subsection{Időigény}

Ha a tömb már rendezett, akkor minden elemnél elég egy összehasonlítás, ezért a legjobb eset:

\[
T(n)=\Theta(n).
\]

Ha a tömb fordított sorrendű, akkor minden új kulcsot végig kell mozgatni a rendezett rész elejéig. Ekkor a mozgatások és összehasonlítások száma nagyságrendileg:

\[
1+2+\dots+(n-1)=\frac{n(n-1)}{2}=\Theta(n^2).
\]

Ezért a legrosszabb és az átlagos eset is:

\[
T(n)=\Theta(n^2).
\]

A beszúró rendezés előnye, hogy egyszerű, helyben rendez, stabil, és kis vagy majdnem rendezett adatsorokra jó választás lehet.

\section{Oszd meg és uralkodj elv}

\subsection{Az elv általános menete}

Az \textbf{oszd meg és uralkodj} stratégia három fő lépésből áll:

\begin{enumerate}
  \item \textbf{Felosztás:} a problémát kisebb, hasonló részproblémákra bontjuk.
  \item \textbf{Uralkodás:} a részproblémákat rekurzívan megoldjuk.
  \item \textbf{Egyesítés:} a részmegoldásokat összerakjuk a teljes probléma megoldásává.
\end{enumerate}

Rendezéseknél ezt az elvet legszebben az összefésülő rendezés és a gyorsrendezés mutatja. Az összefésülő rendezésnél a nehéz lépés az egyesítés, a gyorsrendezésnél pedig a felosztás.

\begin{center}
\begin{tikzpicture}[node distance=9mm]
  \node[flow, minimum width=31mm] (problem) {Teljes tömb\\$n$ elem};
  \node[flow, below left=11mm and 14mm of problem, minimum width=29mm] (left) {Bal rész\\$n/2$};
  \node[flow, below right=11mm and 14mm of problem, minimum width=29mm] (right) {Jobb rész\\$n/2$};
  \node[flow, below=18mm of problem, minimum width=44mm] (merge) {Összefésülés vagy\\particionálás utáni rekurzió};
  \node[flow, below=9mm of merge, minimum width=40mm] (result) {Rendezett tömb};
  \draw[arrow] (problem) -- node[left, font=\small] {felosztás} (left);
  \draw[arrow] (problem) -- node[right, font=\small] {felosztás} (right);
  \draw[arrow] (left) -- (merge);
  \draw[arrow] (right) -- (merge);
  \draw[arrow] (merge) -- node[right, font=\small] {egyesítés / nincs külön egyesítés} (result);
\end{tikzpicture}
\end{center}


\section{Összefésülő rendezés}

\subsection{Alapelv}

Az \textbf{összefésülő rendezés} a tömböt két részre osztja, mindkét részt rekurzívan rendezi, majd a két rendezett részt lineáris időben összefésüli.

Az összefésülés lényege: két rendezett sorozat első még fel nem dolgozott elemét hasonlítjuk össze, és mindig a kisebbet írjuk a kimenet következő helyére. Ha az egyik sorozat elfogy, a másik maradékát egyszerűen átmásoljuk.

\subsection{Pszeudókód}

\zvpseudocode[caption={Összefésülő rendezés},label={lst:zv05_osszefesulo_rendezes}]{src/code/zv05_osszefesulo_rendezes.pseudo}

\zvpseudocode[caption={Két rendezett résztömb összefésülése},label={lst:zv05_osszefesul}]{src/code/zv05_osszefesul.pseudo}

\subsection{Időigény}

Az összefésülő rendezés rekurziós formulája:

\[
T(n)=2T\left(\frac{n}{2}\right)+\Theta(n).
\]

Az $\Theta(n)$ tag az összefésülés költsége. A rekurziós fa minden szintjén összesen $\Theta(n)$ munka történik, a szintek száma pedig $\Theta(\log n)$, ezért:

\[
T(n)=\Theta(n\log n).
\]

Az összefésülő rendezés stabil lehet, ha egyenlőség esetén a bal oldali rész eleme kerül előbb a kimenetbe. Hátránya, hogy általában $\Theta(n)$ segédtárat igényel.

\section{Gyorsrendezés}

\subsection{Alapelv}

A \textbf{gyorsrendezés} is oszd meg és uralkodj típusú algoritmus, de másképpen működik, mint az összefésülő rendezés. Először kiválaszt egy \textbf{pivot} elemet, majd a tömböt úgy osztja két részre, hogy az egyik oldalra kerüljenek a pivotnál nem nagyobb, a másik oldalra pedig a pivotnál nagyobb elemek. Ezután a két részt rekurzívan rendezi.

A gyorsrendezésnél az egyesítés már nem igényel külön munkát, mert a felosztás után a két rész egymáshoz képest jó sorrendben van.

\begin{center}
\begin{tikzpicture}[x=0.72cm,y=0.72cm]
  \foreach \i/\val in {0/3,1/1,2/4,3/2,4/5,5/8,6/7,7/6} {
    \draw[fill=black!5] (\i,0) rectangle +(1,0.8);
    \node at (\i+0.5,0.4) {$\val$};
  }
  \draw[decorate, decoration={brace, mirror}, thick] (0,-0.12) -- (4,-0.12);
  \node at (2,-0.55) {$\leq$ pivot};
  \draw[decorate, decoration={brace, mirror}, thick] (5,-0.12) -- (8,-0.12);
  \node at (6.5,-0.55) {$>$ pivot};
  \draw[fill=warnbg, very thick] (4,0) rectangle +(1,0.8);
  \node at (4.5,0.4) {$5$};
  \node[align=center] at (4.5,1.35) {pivot a végleges\\helyén};
  \draw[->, thick] (4.5,1.05) -- (4.5,0.85);
\end{tikzpicture}
\end{center}


\subsection{Pszeudókód}

\zvpseudocode[caption={Gyorsrendezés},label={lst:zv05_gyorsrendezes}]{src/code/zv05_gyorsrendezes.pseudo}

\zvpseudocode[caption={Particionálás gyorsrendezéshez},label={lst:zv05_particional}]{src/code/zv05_particional.pseudo}

\subsection{Időigény}

A particionálás egy adott részintervallumon lineáris időt igényel. A teljes időigény attól függ, mennyire kiegyensúlyozottak a felosztások.

Legjobb esetben a pivot mindig közel felezi a tömböt:

\[
T(n)=2T\left(\frac{n}{2}\right)+\Theta(n)=\Theta(n\log n).
\]

Legrosszabb esetben a pivot mindig a legkisebb vagy a legnagyobb elem, ezért az egyik rész $0$, a másik $n-1$ elemű:

\[
T(n)=T(n-1)+\Theta(n)=\Theta(n^2).
\]

Átlagosan, jó pivotválasztással vagy véletlenített pivotválasztással:

\[
T(n)=\Theta(n\log n).
\]

A gyorsrendezés a gyakorlatban nagyon gyors lehet, mert helyben dolgozik és jók a konstansai, de általában nem stabil, és rossz pivotválasztással négyzetes időigényűvé válhat.

\section{Buborék rendezés}

\subsection{Alapelv}

A \textbf{buborék rendezés} az egymás mellett álló elemeket hasonlítja össze, és ha rossz sorrendben vannak, felcseréli őket. Egy menet után a legnagyobb elem a rendezendő rész végére kerül. Ezt addig ismétli, amíg nincs több csere.

\zvpseudocode[caption={Buborékrendezés},label={lst:zv05_buborekrendezes}]{src/code/zv05_buborekrendezes.pseudo}

A buborék rendezés előnye, hogy egyszerű és helyben rendez. Ha szerepel benne a cserefigyelés, akkor már rendezett bemeneten $\Theta(n)$ időben megáll. Legrosszabb esetben viszont sok összehasonlítást és sok cserét végez:

\[
T(n)=\Theta(n^2).
\]

\begin{warningbox}
A buborék rendezést vizsgán érdemes egyszerű algoritmusként ismerni, de gyakorlati alkalmazásra nagy adatsorokon általában nem jó választás, mert sok felesleges cserét végez.
\end{warningbox}

\section{Shell rendezés}

\subsection{Alapelv}

A \textbf{Shell rendezés} a beszúró rendezés továbbfejlesztése. A beszúró rendezés gyenge pontja, hogy egy elem csak szomszédos mozgatásokkal jut el messzire. A Shell rendezés ezért eleinte távoli elemeket hasonlít össze és rendez, majd a távolságot fokozatosan csökkenti.

A távolságokat \textbf{növekményeknek} nevezzük. Például:

\[
\dots, 32,16,8,4,2,1
\]

vagy

\[
\dots, 121,40,13,4,1.
\]

Minden növekményre egy olyan beszúró rendezéshez hasonló lépés fut, amely nem szomszédos, hanem $h$ távolságra lévő elemeket rendez. Mire a növekmény $1$ lesz, a tömb már részben rendezett, ezért az utolsó beszúró rendezés gyorsabb lehet.

\subsection{Pszeudókód}

\zvpseudocode[caption={Shell rendezés},label={lst:zv05_shell_rendezes}]{src/code/zv05_shell_rendezes.pseudo}

\subsection{Időigény}

A Shell rendezés időigénye erősen függ a növekménysorozattól. Egyszerű vizsgaszintű megfogalmazásban:

\[
T(n) \approx \Theta(n^{1.5})
\]

alkalmas növekményválasztás mellett, de ez nem univerzális törvény minden növekménysorozatra. A lényeg az, hogy a Shell rendezés a tiszta $\Theta(n^2)$ beszúró rendezésnél sok esetben gyorsabb, de nem éri el általánosan az összefésülő rendezés garantált $\Theta(n\log n)$ idejét.

\section{Minimum kiválasztásos rendezés}

\subsection{Alapelv}

A \textbf{minimum kiválasztásos rendezés} minden lépésben megkeresi a rendezendő rész legkisebb elemét, és azt a rész elejére cseréli. Az első menet után a legkisebb elem a helyén van, a második menet után a második legkisebb is, és így tovább.

\zvpseudocode[caption={Minimum kiválasztásos rendezés},label={lst:zv05_minimum_kivalasztasos_rendezes}]{src/code/zv05_minimum_kivalasztasos_rendezes.pseudo}

Az összehasonlítások száma bemenettől függetlenül:

\[
(n-1)+(n-2)+\dots+1=\frac{n(n-1)}{2}=\Theta(n^2).
\]

A cserék száma viszont legfeljebb $n-1$, ezért akkor lehet érdekes, ha az összehasonlítás olcsó, de a mozgatás drága. Általában helyben rendez, de nem stabil, mert egy távoli minimum előrecserélése felboríthatja az azonos kulcsú elemek sorrendjét.

\section{Négyzetes rendezés}

A tételben szereplő \textbf{négyzetes rendezés} nem azt jelenti, hogy az időigénye $n^2$, hanem azt, hogy a módszer a bemenetet közel $\sqrt n$ darab, közel $\sqrt n$ méretű részre osztja. Az alapgondolat:

\begin{enumerate}
  \item az $n$ elemet több alcsoportba osztjuk;
  \item minden alcsoportból kiemeljük a legkisebb elemet;
  \item ezekből főcsoportot képezünk;
  \item mindig a főcsoport minimumát írjuk az eredménybe;
  \item abból az alcsoportból, ahonnan a minimum származott, új elemet emelünk be.
\end{enumerate}

\subsection{Pszeudókód}

\zvpseudocode[caption={Négyzetes rendezés},label={lst:zv05_negyzetes_rendezes}]{src/code/zv05_negyzetes_rendezes.pseudo}

Ha $\sqrt n$ alcsoport van, és mindegyik közel $\sqrt n$ méretű, akkor a módszer összehasonlításszáma:

\[
T(n)=\Theta(n\sqrt n)=\Theta(n^{1.5}).
\]

Továbbfejlesztett, több szintű csoportosítással ez javítható például $\Theta(n^{4/3})$ nagyságrendre. A vizsgán itt a lényeg az, hogy a módszer a minimumkiválasztás ötletét csoportosítással gyorsítja, de az általános összehasonlító alsó korlát miatt nem tud valódi lineáris összehasonlító rendezéssé válni.

\section{Lineáris idejű rendezések}

\subsection{Miért lehetnek lineárisak?}

A leszámláló és számjegyes rendezés nem pusztán összehasonlításokra épül. Ezek kulcsszerkezeti feltételeket használnak ki, például azt, hogy a kulcsok kis tartományba eső egészek, vagy fix hosszúságú számjegysorozatok.

Ezért nem mondanak ellent az összehasonlító rendezések $\Omega(n\log n)$ alsó korlátjának.

\begin{center}
\begin{tikzpicture}[node distance=8mm]
  \node[flow, minimum width=42mm] (input) {Bemenet\\$A[1..n]$};
  \node[flow, right=11mm of input, minimum width=42mm] (count) {Számlálás\\$C[1..k]$};
  \node[flow, right=11mm of count, minimum width=42mm] (prefix) {Prefixösszegek\\$C[i]=\#\{x\leq i\}$};
  \node[flow, below=11mm of prefix, minimum width=42mm] (place) {Elemek végleges\\pozícióba írása};
  \node[flow, left=11mm of place, minimum width=42mm] (output) {Rendezett kimenet\\$B[1..n]$};
  \draw[arrow] (input) -- (count);
  \draw[arrow] (count) -- (prefix);
  \draw[arrow] (prefix) -- (place);
  \draw[arrow] (place) -- (output);
\end{tikzpicture}
\end{center}


\subsection{Leszámláló rendezés}

A \textbf{leszámláló rendezés} bemenete olyan egész kulcsokból áll, amelyek $1$ és $k$ közé esnek. Három tömböt használ:

\begin{itemize}
  \item $A[1..n]$: bemeneti tömb;
  \item $B[1..n]$: kimeneti tömb;
  \item $C[1..k]$: számláló tömb.
\end{itemize}

A módszer lépései:

\begin{enumerate}
  \item lenullázza a $C$ tömböt;
  \item megszámolja, hogy melyik kulcs hányszor szerepel;
  \item prefixösszegeket képez, így $C[i]$ megadja, hány elem $\leq i$;
  \item hátulról előre bejárja a bemenetet, és minden elemet a végleges helyére ír a $B$ tömbben.
\end{enumerate}

\zvpseudocode[caption={Leszámláló rendezés},label={lst:zv05_leszamlalo_rendezes}]{src/code/zv05_leszamlalo_rendezes.pseudo}

Az időigény:

\[
T(n)=\Theta(n+k).
\]

Ha $k=\Theta(n)$, akkor:

\[
T(n)=\Theta(n).
\]

A leszámláló rendezés stabil, ha a bemenetet hátulról előre járjuk be a kimenetbe íráskor. Ez különösen fontos a számjegyes rendezésnél, mert ott a stabil részrendezések egymásra épülnek.

\subsection{Számjegyes rendezés}

A \textbf{számjegyes rendezés}, más néven radix rendezés, fix hosszúságú kulcsokat rendez több menetben. Például számokat számjegyenként vagy sztringeket karakterenként dolgozhat fel.

A leggyakoribb változat a legkisebb helyiértéktől halad a legnagyobb felé. Minden pozíció szerint stabil rendezést végzünk.

\zvpseudocode[caption={Számjegyes rendezés},label={lst:zv05_szamjegyes_rendezes}]{src/code/zv05_szamjegyes_rendezes.pseudo}

Ha $d$ a számjegyek száma, $k$ pedig egy számjegy lehetséges értékeinek száma, akkor az időigény:

\[
T(n)=\Theta(d(n+k)).
\]

Ha $d$ és $k$ az $n$-hez képest kicsi vagy konstansnak tekinthető, akkor a számjegyes rendezés lineáris jellegű.

\begin{center}
\begingroup
\setlength{\tabcolsep}{4pt}
\begin{tabularx}{\textwidth}{L{0.22\textwidth}X X L{0.18\textwidth}}
\toprule
Módszer & Kulcsfeltétel & Miért lehet gyors? & Időigény \\
\midrule
Leszámláló rendezés & A kulcs egész szám, és $1..k$ közötti tartományba esik. & Nem elempárokat hasonlít, hanem előfordulásszámokat és prefixösszegeket használ. & $\Theta(n+k)$ \\
Számjegyes rendezés & A kulcs $d$ darab pozícióra vagy számjegyre bontható, pozíciónként $k$ lehetséges értékkel. & Több stabil, egyszerű részrendezést végez egymás után. & $\Theta(d(n+k))$ \\
Edényrendezés & A bemenet eloszlásáról is feltételezünk valamit, például egyenletes eloszlást. & Az elemeket edényekbe szórja, és az edényeken belül kicsi rendezéseket végez. & várható $\Theta(n)$ \\
\bottomrule
\end{tabularx}
\endgroup

\end{center}

\section{Az összehasonlító rendezések alsó korlátja}

\subsection{Döntési fa modell}

Egy összehasonlító rendezés működése modellezhető döntési fával. A fa belső csúcsai összehasonlítások, például:

\[
A[i] \leq A[j]?
\]

A bal vagy jobb ág annak felel meg, hogy az összehasonlítás eredménye igaz vagy hamis. A fa levelei a lehetséges rendezett permutációknak felelnek meg.

\begin{center}
\begin{tikzpicture}[level distance=12mm,
  level 1/.style={sibling distance=42mm},
  level 2/.style={sibling distance=22mm},
  every node/.style={align=center}]
  \node[pred] {$a_1\leq a_2?$}
    child { node[pred] {$a_2\leq a_3?$}
      child { node[flow, minimum width=22mm] {$a_1a_2a_3$} }
      child { node[flow, minimum width=22mm] {$a_1a_3a_2$} }
    }
    child { node[pred] {$a_1\leq a_3?$}
      child { node[flow, minimum width=22mm] {$a_2a_1a_3$} }
      child { node[flow, minimum width=22mm] {$\dots$} }
    };
  \node[align=center, text width=0.86\textwidth] at (0,-5.2) {A leveleknek a lehetséges permutációkat kell megkülönböztetniük:\\legalább $n!$ levél szükséges.};
\end{tikzpicture}
\end{center}


\subsection{Az alsó korlát gondolata}

$n$ különböző elemnek $n!$ lehetséges sorrendje van. Egy helyes összehasonlító rendezésnek minden lehetséges bemeneti permutációt meg kell tudnia különböztetni, ezért a döntési fának legalább $n!$ levele kell legyen.

Egy $h$ magasságú bináris fa legfeljebb $2^h$ levelet tartalmaz, ezért:

\[
2^h \geq n!.
\]

Logaritmust véve:

\[
h \geq \log_2(n!).
\]

Stirling-formulával:

\[
\log_2(n!)=\Omega(n\log n).
\]

Tehát bármely összehasonlító rendezés legrosszabb esetben legalább

\[
\Omega(n\log n)
\]

összehasonlítást igényel.

\begin{warningbox}
Ez az alsó korlát csak az összehasonlító rendezésekre vonatkozik. A leszámláló rendezés és a számjegyes rendezés azért lehet lineáris, mert a kulcsokról többet feltételez, mint pusztán azt, hogy összehasonlíthatók.
\end{warningbox}

\section{Összefoglalás vizsgára}

\begin{center}
\begingroup
\setlength{\tabcolsep}{3.5pt}
\begin{longtable}{L{0.20\textwidth}L{0.16\textwidth}L{0.16\textwidth}L{0.16\textwidth}L{0.14\textwidth}L{0.08\textwidth}}
\toprule
Algoritmus & Legjobb eset & Átlagos eset & Legrosszabb eset & Segédtár & Stabil? \\
\midrule
\endfirsthead
\toprule
Algoritmus & Legjobb eset & Átlagos eset & Legrosszabb eset & Segédtár & Stabil? \\
\midrule
\endhead
Beszúró rendezés & $\Theta(n)$ & $\Theta(n^2)$ & $\Theta(n^2)$ & $\Theta(1)$ & igen \\
Összefésülő rendezés & $\Theta(n\log n)$ & $\Theta(n\log n)$ & $\Theta(n\log n)$ & $\Theta(n)$ & igen, ha így írjuk \\
Gyorsrendezés & $\Theta(n\log n)$ & $\Theta(n\log n)$ & $\Theta(n^2)$ & rekurziós verem & nem \\
Buborék rendezés & $\Theta(n)$ & $\Theta(n^2)$ & $\Theta(n^2)$ & $\Theta(1)$ & igen \\
Shell rendezés & növekményfüggő & növekményfüggő & növekményfüggő & $\Theta(1)$ & nem általában \\
Minimum kiválasztásos & $\Theta(n^2)$ & $\Theta(n^2)$ & $\Theta(n^2)$ & $\Theta(1)$ & nem általában \\
Négyzetes rendezés & kb. $\Theta(n^{1.5})$ & kb. $\Theta(n^{1.5})$ & kb. $\Theta(n^{1.5})$ & csoportfüggő & változó \\
Leszámláló rendezés & $\Theta(n+k)$ & $\Theta(n+k)$ & $\Theta(n+k)$ & $\Theta(n+k)$ & igen \\
Számjegyes rendezés & $\Theta(d(n+k))$ & $\Theta(d(n+k))$ & $\Theta(d(n+k))$ & segédrendezéstől függ & igen, ha a részrendezés stabil \\
\bottomrule
\end{longtable}
\endgroup

\end{center}

Rövid, vizsgán jól használható összehasonlítás:

\begin{itemize}
  \item A \textbf{beszúró rendezés} egyszerű, stabil, helyben rendez, majdnem rendezett bemeneten gyors, de általánosan $\Theta(n^2)$.
  \item Az \textbf{összefésülő rendezés} garantáltan $\Theta(n\log n)$, stabil lehet, de $\Theta(n)$ segédtár kell hozzá.
  \item A \textbf{gyorsrendezés} átlagosan $\Theta(n\log n)$ és gyakorlatban gyors, de legrosszabb esetben $\Theta(n^2)$.
  \item A \textbf{buborék rendezés} egyszerű, de sok cserét végez, ezért inkább elméleti alapalgoritmus.
  \item A \textbf{Shell rendezés} a beszúró rendezést gyorsítja távoli cserékkel, ideje a növekménysorozattól függ.
  \item A \textbf{minimum kiválasztásos rendezés} mindig $\Theta(n^2)$ összehasonlítást végez, de kevés cserét használ.
  \item A \textbf{négyzetes rendezés} csoportosítással gyorsítja a minimumkiválasztás gondolatát, tipikusan $\Theta(n^{1.5})$ nagyságrenddel.
  \item A \textbf{leszámláló rendezés} $\Theta(n+k)$ idejű, ha a kulcsok kis egész tartományból jönnek.
  \item A \textbf{számjegyes rendezés} $\Theta(d(n+k))$ idejű, és stabil részrendezést igényel.
  \item Az \textbf{összehasonlító rendezések alsó korlátja} $\Omega(n\log n)$, ezért általános összehasonlító módszerrel ennél aszimptotikusan jobb legrosszabb eseti rendezés nem lehetséges.
\end{itemize}
